\begin{lemma}
\label{editorial-source-lemma-charpoly-module}
Let $R$ be a ring.
Let $M$ be a finite $R$-module.
Let $\varphi : M \to M$ be an endomorphism.
Then there exists a monic polynomial $P \in R[T]$ such that
$P(\varphi) = 0$ as an endomorphism of $M$.
\end{lemma}

\begin{proof}
Choose a surjective $R$-module map $R^{\oplus n} \to M$, given by
$(a_1, \ldots, a_n) \mapsto \sum a_ix_i$ for some generators $x_i \in M$.
Choose $(a_{i1}, \ldots, a_{in}) \in R^{\oplus n}$ such that
$\varphi(x_i) = \sum a_{ij} x_j$. In other words the diagram
$$
\xymatrix{
R^{\oplus n} \ar[d]_{A^{\mathrm{t}}} \ar[r] & M \ar[d]^\varphi \\
R^{\oplus n} \ar[r] & M
}
$$
is commutative where $A = (a_{ij})$ and $A^{\mathrm{t}}$
denotes its transpose: the lift sends the $i$th basis vector to
$\sum_j a_{ij}e_j$, which maps to $\varphi(x_i)$.
By
Lemma \ref{algebra-lemma-charpoly}
there exists a monic polynomial $P$ such that $P(A) = 0$.
Taking transposes gives $P(A^{\mathrm{t}})=P(A)^{\mathrm{t}}=0$,
since the coefficient matrices are scalar and transpose preserves
powers of this one matrix.
Then it follows that $P(\varphi) = 0$.
\end{proof}